\subsection{Is the Route Optimal?}\label{sec:optimality}
Divide and Conquer gives no proof of optimality, so we checked the result separately. With only
\NumVertices{} vertices, the exact Held--Karp algorithm~\cite{heldkarp1962} (Algorithm~2.2 of the
paper) can be run in well under a second. It gives the provably optimal tour
\begin{center}
\OptRoute{}
\end{center}
of length \textbf{\OptTotal{} km}. Our Divide and Conquer route is therefore \textbf{not optimal}: it is
\DCTotal{} km, i.e.\ \DCGap\% longer than the optimum. It is an \emph{approximate} solution.

Because the algorithm is randomised, we also repeated it with \RunCount{} different seeds, as the paper does
in its Table~2. Table~\ref{tab:experiments} summarises the results.

\begin{table}[H]
\centering
\caption{Tour length (km) over \RunCount{} seeds, compared with the exact optimum. ``Optimal'' counts the runs
that reached the optimum.}\label{tab:experiments}
\vspace{4pt}
\small
% AUTO-GENERATED by code/main.py --export-latex. Do not edit; re-run the export instead.
\begin{tabular}{l r r r r r r r}
\toprule
\textbf{Method} & \textbf{Best} & \textbf{Worst} & \textbf{Mean} & \textbf{Std.} & \textbf{Mean gap} & \textbf{Optimal} & \textbf{ms/run} \\
\midrule
Shortest bridge (paper) & 13.61 & 14.26 & 13.88 & 0.25 & 2.01\% & 6/100 & 0.405 \\
Best splice (variant) & 13.61 & 14.13 & 13.91 & 0.24 & 2.20\% & 17/100 & 0.342 \\
Held--Karp (exact) & 13.61 & 13.61 & 13.61 & 0.00 & 0.00\% & -- & -- \\
\bottomrule
\end{tabular}

\end{table}

With the paper's shortest-bridge merge, the tour length ranged from \BridgeBest{} to \BridgeWorst{}~km
(mean \BridgeMean{}~km, \BridgeMeanGap\% above optimal). Only \BridgeOptHits{} of the \RunCount{} runs
happened to produce an optimal tour, and the algorithm itself cannot recognise when this happens; only
the separate Held--Karp check shows it. This confirms that the method is a heuristic: on its own it
gives a good tour, not a proven optimal one.
